% !TEX program = lualatex
\documentclass[aspectratio=43,professionalfonts,handout]{beamer}
\usepackage{beamer_template}

\newcommand{\LectureNum}{6}
\newcommand{\LectureTitle}{たたみ込み}
\newcommand{\TextPages}{p.\,65-67}
\newcommand{\CourseName}{応用数学}
\renewcommand{\TermName}{前期}

\begin{document}

\begin{frame}{\CourseName\quad 第\LectureNum 回「\LectureTitle」}
  {\small \TermName\quad \TeacherName\quad ／\quad 教科書 \TextPages}

  \textbf{目標}：たたみ込みの定義と性質を理解し、たたみ込みのラプラス変換を用いた計算ができる

\end{frame}

\begin{frame}{たたみ込みの定義}
  \begin{block}{}
  \begin{itemize}
    \item $(f*g)(t) = \int_0^t f(\tau)g(t-\tau)d\tau$
    \item 交換法則：$f*g = g*f$
    \item グラフのイメージ
  \end{itemize}
  \end{block}
\end{frame}

\begin{frame}{たたみ込みのラプラス変換}
  \begin{block}{}
  \begin{itemize}
    \item $\mathcal{L}[f*g] = \mathcal{L}[f]$ ・ $\mathcal{L}[g] = F(s)G(s)$
    \item 証明の概要
    \item 逆変換への応用：$\mathcal{L}^{-1}[F(s)G(s)] = f*g$
  \end{itemize}
  \end{block}
\end{frame}

\begin{frame}{積分方程式への応用}
  \begin{block}{}
  \begin{itemize}
    \item 積分記号内に未知関数を含む方程式
    \item 例題8：左辺は x*sin のたたみ込み
    \item $\mathcal{L}$ 変換 $ \to $ 代数方程式 $ \to $ 逆変換
  \end{itemize}
  \end{block}
\end{frame}

\begin{frame}{演習}
  \begin{exampleblock}{自分で解いてみよう}
  \begin{itemize}
    \item 問5〜9より選題
  \end{itemize}
  \end{exampleblock}
\end{frame}

\begin{frame}{まとめ}
  \begin{itemize}
    \item たたみ込みの定義と性質
    \item $\mathcal{L}[f*g] = F(s)G(s)$
    \item 次回：線形システムとデルタ関数
  \end{itemize}
\end{frame}

\end{document}
